\chapter{Rezoluție și Prolog}

\section{Termeni compusi in Prolog}

Pana acum in prolog, am lucrat doar cu termeni simpli (variabile, constante, atomi), dar in prolog mai
exista si conceptul de \textbf{termeni compusi}.

De exemplu putem codifica formula logica $\varphi = (v_0 \lor v_1) \rightarrow v_2$ astfel:

\texttt{implica(sau(v0, v1), v2)}

In acest exemplu, \texttt{implica} \textbf{nu este un fapt} (nu avem punct la final) ci un termen compus.

\vspace{\baselineskip}

Ca sa intelegem cum se lucreaza cu ele, putem lua urmatorul exemplu unde vrem sa aflam
toate variabilele unei expresii logice (pentru $\varphi$ definit anterior variabilele ar fi \texttt{[v0, v1, v2]}).
Pentru simplitate vom afisa si duplicate.

\begin{lstlisting}
vars(A, [A]) :- atom(A).

vars(not(A), R) :- vars(A, R).

vars(si(A, B), R) :- 
    vars(A, R1),
    vars(B, R2),
    append(R1, R2, R).

vars(sau(A, B), R) :- 
    vars(A, R1),
    vars(B, R2),
    append(R1, R2, R).

vars(implica(A, B), R) :- 
    vars(A, R1),
    vars(B, R2),
    append(R1, R2, R).
\end{lstlisting}

\texttt{?- vars(implica(sau(v0, v1), v2), R).}

\texttt{R = [v0, v1, v2].}

\texttt{?- vars(sau(si(v1, v2), si(v2, v4)), R).}

\texttt{R = [v1, v2, v2, v4].}

\newpage

\subsection*{Predicate utile}

\textbf{Predicatul \texttt{functor(termen compus, functor, aritate)}} extrage functorul si aritatea
termenului compus:

\texttt{?- functor(data(22, aprilie, 2026), F, A).}

\texttt{F = data,}

\texttt{A = 3.}

\vspace{\baselineskip}

\textbf{Predicatul \texttt{arg(i, termen compus, rez)}} extrage termenul de la pozitia i din termenul 
compus (atentie, este indexat de la 1):

\texttt{?- arg(2, data(22, aprilie, 2026), Luna).}

\texttt{Luna = aprilie.}

\vspace{\baselineskip}

\textbf{Operatorul \texttt{=..}} transforma un termen compus intr-o lista si invers:

\texttt{?- data(22, aprilie, 2026) =.. Lista.}

\texttt{Lista = [data, 22, aprilie, 2026].}

\texttt{?- Data =.. [data, 22, aprilie, 2026].}

\texttt{Data = data(22, aprilie, 2026).}


\subsection*{Exercitii propuse:}

\subsubsection*{Exercitiul 1}
Definiti predicatul \texttt{suma\_noduri(Arbore, Suma)} care calculeaza suma nodurilor
unui arbore binar. Folositi termenul compus \texttt{nod(Valoare, ArbStang, ArbDrept)}.

\texttt{?- suma\_noduri(nod(5, nod(3, nod(1, nil, nil), nod(4, nil, nil)), nod(8, nil, nil)), R).}

\texttt{R = 21.}

\subsubsection*{Exercitiul 2}
Definiti predicatul \texttt{substituie(ExpresieInitiala, CeInlocuiesc, CuCeInlocuiesc, ExpresieFinala)} care
sparge recursiv o structura si inlocuieste aparitiile unui subtermen specific cu altul.

\texttt{?- substituie(implica(v1, si(v1, v3)), v1, not(v2), R).}

\texttt{R = implica(not(v2), si(not(v2), v3)).}

\subsubsection*{Exercitiul 3}
Definiti predicatul \texttt{evaluare(Expresie, ListaAdevar, Rezultat)} care primeste
o expresie logica si o lista de variabile care sunt considerate adevarate.

Rulati si confirmati:

\texttt{?- evaluare(si(v1, sau(v2, v3)), [v1, v3], R).}

\texttt{R = true.}

\texttt{?- evaluare(imp(v3, si(v1, sau(v2, v3))), [v1], R).}

\texttt{R = false.}




\newpage
\begin{ex}
Să se demonstreze următoarele teoreme prin metoda deducției naturale:
\begin{enumerate}
    \item[(i)] $(\varphi \rightarrow (\psi \rightarrow \chi)) \rightarrow (\psi \rightarrow (\varphi \rightarrow \chi))$;
    \item[(ii)] $(\varphi \rightarrow \psi) \rightarrow (\varphi \rightarrow (\chi \rightarrow \psi))$.
\end{enumerate}
\end{ex}

\begin{solutie}
\vspace{0.3cm}
\noindent \textbf{Pentru formula (i):} $\vdash (\varphi \rightarrow (\psi \rightarrow \chi)) \rightarrow (\psi \rightarrow (\varphi \rightarrow \chi))$

Pentru a urmări demonstrația clar, prima coloană va reprezenta starea curentă (secvența deductivă), iar a doua coloană va indica justificarea (regula aplicată).

\begin{center}
\renewcommand{\arraystretch}{1.3} % Spațiere puțin mai mare între rânduri pentru lizibilitate
\begin{tabular}{r l l}
1) & $\{\varphi \rightarrow (\psi \rightarrow \chi), \psi, \varphi\} \vdash \varphi \rightarrow (\psi \rightarrow \chi)$ & P1.53(ii) \\
2) & $\{\varphi \rightarrow (\psi \rightarrow \chi), \psi, \varphi\} \vdash \varphi$ & P1.53(ii) \\
3) & $\{\varphi \rightarrow (\psi \rightarrow \chi), \psi, \varphi\} \vdash \psi \rightarrow \chi$ & MP: 1), 2) \\
4) & $\{\varphi \rightarrow (\psi \rightarrow \chi), \psi, \varphi\} \vdash \psi$ & P1.53(ii) \\
5) & $\{\varphi \rightarrow (\psi \rightarrow \chi), \psi, \varphi\} \vdash \chi$ & MP: 3), 4) \\
6) & $\{\varphi \rightarrow (\psi \rightarrow \chi), \psi\} \vdash \varphi \rightarrow \chi$ & Th. Ded. din 5) \\
7) & $\{\varphi \rightarrow (\psi \rightarrow \chi)\} \vdash \psi \rightarrow (\varphi \rightarrow \chi)$ & Th. Ded. din 6) \\
8) & $\vdash (\varphi \rightarrow (\psi \rightarrow \chi)) \rightarrow (\psi \rightarrow (\varphi \rightarrow \chi))$ & Th. Ded. din 7) \\
\end{tabular}
\end{center}

\vspace{0.5cm}
\noindent \textbf{Pentru formula (ii):} $\vdash (\varphi \rightarrow \psi) \rightarrow (\varphi \rightarrow (\chi \rightarrow \psi))$

\begin{center}
\renewcommand{\arraystretch}{1.3}
\begin{tabular}{r l l}
1) & $\{\varphi \rightarrow \psi, \varphi\} \vdash \varphi \rightarrow \psi$ & P1.53(ii) \\
2) & $\{\varphi \rightarrow \psi, \varphi\} \vdash \varphi$ & P1.53(ii) \\
3) & $\{\varphi \rightarrow \psi, \varphi\} \vdash \psi$ & MP: 1), 2) \\
4) & $\{\varphi \rightarrow \psi, \varphi\} \vdash \psi \rightarrow (\chi \rightarrow \psi)$ & A1 \\
5) & $\{\varphi \rightarrow \psi, \varphi\} \vdash \chi \rightarrow \psi$ & MP: 3), 4) \\
6) & $\{\varphi \rightarrow \psi\} \vdash \varphi \rightarrow (\chi \rightarrow \psi)$ & Th. Ded. din 5) \\
7) & $\vdash (\varphi \rightarrow \psi) \rightarrow (\varphi \rightarrow (\chi \rightarrow \psi))$ & Th. Ded. din 6) \\
\end{tabular}
\end{center}

\end{solutie}